\documentclass[11pt,a4paper]{article}

% ─── packages ───
\usepackage[utf8]{inputenc}
\usepackage[T1]{fontenc}
\usepackage{lmodern}
\usepackage{amsmath,amssymb}
\usepackage{microtype}
\usepackage{booktabs}
\usepackage{hyperref}
\usepackage{graphicx}
\usepackage{natbib}
\usepackage[margin=1in]{geometry}
\usepackage{xcolor}
\usepackage{float}
\usepackage{caption}

\hypersetup{
  colorlinks=true,
  linkcolor=blue!60!black,
  citecolor=blue!60!black,
  urlcolor=blue!60!black
}

\title{Improving Curve Fitting for AI Capability Time Horizons:\\
       Three Drop-In Improvements to METR's Methodology}
\author{Gabor Gulyas}
\date{March 2026}

\begin{document}
\maketitle

\begin{abstract}
METR's Time Horizon benchmark is the most widely cited tracker of AI agent capability growth, with its headline finding---a capability doubling time of approximately four months---informing AI safety policy.
We propose three concrete, drop-in improvements to METR's curve-fitting pipeline: (1)~replacing logistic regression with isotonic regression to avoid the symmetric-sigmoid assumption, (2)~replacing the $p_{50}$ threshold-crossing summary metric with a geometric-mean integral metric $G[T]$ that uses the entire fitted curve, and (3)~using $m$-out-of-$n$ bootstrap instead of standard bootstrap for theoretical consistency with shape-constrained estimators.
We validate each improvement with 5-fold cross-validation on 18 frontier AI models.
Isotonic regression improves cross-validated Brier score from $0.1212$ to $0.1115$ and cross-validated log-loss from $0.376$ to $0.368$.
The $G[T]$ metric yields a doubling-time estimate of $128$--$130$ days with $R^2 \geq 0.916$ across both methods, compared to $118$--$121$ days with $R^2$ ranging from $0.811$ to $0.916$ under $p_{50}$.
The good news: the headline conclusion---AI agent capabilities are doubling roughly every four months---is robust to all methodology choices we tested.
The improvements make the estimate more defensible, not radically different.

\medskip\noindent
\textbf{Interactive companion:} All per-model fits, trend charts, and comparison tables are available at an interactive viewer:
\url{https://guyko81.github.io/metr-ablation/} \hfill {\small(GitHub Pages)}

\medskip\noindent
\textbf{Reproducibility:} Full code and data at \url{https://github.com/guyko81/metr-ablation}.
\end{abstract}

% ══════════════════════════════════════════════════════════════
\section{Introduction}
\label{sec:intro}

METR's Time Horizon benchmark~\citep{metr2025} is the most important systematic effort to track how the capabilities of AI agents evolve over time.
By measuring each model's success rate on a battery of real-world tasks spanning a wide range of human completion times (from minutes to days), METR estimates a \emph{capability time horizon}---the task duration at which a model reaches a given success rate---and tracks how this horizon grows with each new model release.

The headline result, that AI agent capabilities are doubling approximately every four months, is increasingly cited in policy discussions, safety assessments, and capability forecasting.
Because this number matters, the statistical methodology behind it should be as robust as possible.

We identify three specific methodology choices in METR's pipeline where principled alternatives exist, and we show that adopting them yields better-calibrated fits and more stable trend estimates.
Crucially, none of our improvements change the headline conclusion.
The doubling time moves from $121$ days (logistic + $p_{50}$) to $130$ days (isotonic + $G[T]$)---still approximately four months---but the estimate becomes more defensible against methodological critique.

Our three improvements are:
\begin{enumerate}
  \item \textbf{Curve fitting:} Replace logistic regression with isotonic regression to avoid the symmetric-sigmoid assumption.
  \item \textbf{Summary metric:} Replace the $p_{50}$ threshold-crossing point with $G[T]$, an integral metric that summarizes the entire fitted curve.
  \item \textbf{Bootstrap method:} Use $m$-out-of-$n$ bootstrap (with $m = n^{2/3}$) instead of standard bootstrap when using isotonic estimators, following established theory for shape-constrained inference.
\end{enumerate}

All code, data, and interactive visualizations are publicly available.\footnote{Interactive viewer: \url{https://guyko81.github.io/metr-ablation/}; repository: \url{https://github.com/guyko81/metr-ablation}}


% ══════════════════════════════════════════════════════════════
\section{Background}
\label{sec:background}

METR's pipeline proceeds in four stages:
\begin{enumerate}
  \item \textbf{Data collection.} Each AI model is evaluated on a battery of tasks with known human completion times, producing binary success/failure outcomes weighted by task difficulty.
  \item \textbf{Curve fitting.} A success-rate curve $p(u)$ is fitted, where $u = \log_2(\text{human minutes})$, mapping task difficulty to estimated success probability.
  \item \textbf{Summary statistic.} A scalar ``capability time horizon'' is extracted from each model's fitted curve---in METR's case, the $p_{50}$ crossing point (the task duration at which $p(u) = 0.5$).
  \item \textbf{Trend regression.} An exponential trend is fitted to the $(\text{release date}, \text{capability horizon})$ pairs across models, yielding the doubling-time estimate.
\end{enumerate}

METR uses logistic regression for step~2 and the $p_{50}$ crossing point for step~3.
We propose improvements to both choices.

Our analysis covers 18 frontier AI models released between March 2023 and December 2025, including the Claude family (3 Opus through Opus~4.5), the GPT family (GPT-4 through GPT-5.2), reasoning models (o1-preview, o1, o3), and Gemini~3~Pro, with approximately 1{,}000--1{,}400 evaluated task runs per model.


% ══════════════════════════════════════════════════════════════
\section{Three Improvements}
\label{sec:improvements}

% ──────────────────────────────────────────────────────────────
\subsection{Curve Fitting: Isotonic Regression}
\label{sec:isotonic}

\paragraph{The issue.}
Logistic regression fits a sigmoid:
\[
  p(u) = \frac{1}{1 + \exp(-(\beta_0 + \beta_1 u))}
\]
This is symmetric around its inflection point---it assumes the curve drops from high success to low success at the same rate in both directions.
In practice, success-versus-difficulty curves are often asymmetric: models may saturate quickly on easy tasks (success near 1 over a broad range) but exhibit a long, slow tail on hard tasks.
A symmetric sigmoid cannot capture this shape, biasing both the fitted curve and any threshold crossing derived from it.

\paragraph{The fix.}
Isotonic regression~\citep{barlow1972,robertson1988} fits a monotone non-increasing function with no parametric assumption.
It preserves only the structural prior that harder tasks should not be easier, and adapts its shape freely to the data.
Raw isotonic regression produces a step function, which is discontinuous; however, the bootstrap confidence bands provide sufficient smoothing for visual interpretation, and the $G[T]$ integral metric (Section~\ref{sec:gt}) averages over the entire curve, making the step-function discontinuities irrelevant for trend estimation.

\paragraph{Evidence.}
Table~\ref{tab:results} summarizes fit quality.
Using 5-fold cross-validation (averaging across all 18 models):

\begin{itemize}
  \item \textbf{Brier score:} Isotonic achieves $0.1115$ vs.\ logistic's $0.1212$ (8.0\% improvement).
  \item \textbf{Log-loss:} Isotonic achieves $0.368$ vs.\ logistic's $0.376$ (2.1\% improvement).
  \item Both improvements hold \emph{out of sample}, confirming they are not artifacts of in-sample overfitting.
\end{itemize}

\paragraph{Overfitting check.}
The in-sample Brier gap ($0.1085$ vs.\ $0.1204$, a 9.9\% gap) is larger than the cross-validated gap ($0.1115$ vs.\ $0.1212$, an 8.0\% gap).
This confirms isotonic regression does overfit slightly more than logistic---as expected for a nonparametric method---but the improvement survives cross-validation, indicating genuine better calibration.


% ──────────────────────────────────────────────────────────────
\subsection{Summary Metric: $G[T]$ Integral}
\label{sec:gt}

\paragraph{The issue.}
The $p_{50}$ crossing point---the task duration $T$ at which $p(\log_2 T) = 0.5$---is a single-point statistic.
It depends entirely on the fitted curve's behavior at one specific success level.
For step-function curves (isotonic regression), $p_{50}$ can jump discontinuously between adjacent steps, introducing high variance.
Even for smooth curves, small perturbations near the 50\% level can shift $p_{50}$ substantially.
Indeed, switching from logistic to isotonic regression changes the per-model $p_{50}$ values dramatically---for example, Claude~Opus~4.5 shifts from $320$ to $119$ minutes---even though the overall doubling time remains similar ($121$ vs.\ $118$ days).
This sensitivity to methodology at the individual-model level, while the trend is preserved, underscores that $p_{50}$ is not a reliable per-model summary.

\paragraph{The fix.}
We define the \emph{geometric mean capability time} $G[T]$ as follows.
Let $q(u) = 1 - p(u)$ serve as a CDF-like function over the log$_2$-difficulty axis $u$.
Compute:
\[
  E[U] = u_{\min} + \int_{u_{\min}}^{u_{\max}} p(u)\, du
\]
where $u_{\min}$ and $u_{\max}$ are the endpoints of the evaluation grid.
Then:
\[
  G[T] = 2^{E[U]}
\]
This integral uses the \emph{entire fitted curve}, averaging over the bulk of the difficulty range where all three methods agree closely.
The exponentiation converts back from log$_2$-space to minutes.

\paragraph{Why it works.}
The integral averages over the full range of task difficulties, washing out local fluctuations.
In regions where the models are easy (success $\approx 1$) or hard (success $\approx 0$), all methods agree.
The methods diverge most in the transition zone around $p = 0.5$, which is exactly the single point $p_{50}$ depends on.
$G[T]$ dilutes this sensitive region with the stable bulk, producing a much more robust summary statistic.

\paragraph{Evidence.}
\begin{itemize}
  \item $G[T]$ doubling time: $128$--$130$ days across both methods ($R^2 = 0.916$--$0.921$).
  \item $p_{50}$ doubling time: $118$--$121$ days ($R^2 = 0.811$--$0.916$).
  \item The $R^2$ spread across methods shrinks from $0.105$ ($p_{50}$) to $0.005$ ($G[T]$).
\end{itemize}

Although the $p_{50}$ doubling time is roughly the same across methods ($\sim$\!$120$ days), the actual $p_{50}$ \emph{levels} change significantly between logistic and isotonic fits for individual models.
$G[T]$ is far more stable at the per-model level because it integrates over the entire curve, capturing the full picture in a single number rather than depending on behavior at one arbitrary threshold.

One could argue that $G[T]$'s stability reflects insensitivity to tail behavior.
For a headline trend metric, this is a feature: the doubling-time estimate should be robust across reasonable methodology choices, and $G[T]$ achieves this.


% ──────────────────────────────────────────────────────────────
\subsection{Bootstrap: $m$-out-of-$n$ Resampling}
\label{sec:bootstrap}

Standard bootstrap draws $n$ samples with replacement from a dataset of size $n$.
For isotonic regression, this is theoretically inconsistent: the isotonic estimator converges at rate $n^{1/3}$ (not $n^{1/2}$ as for smooth parametric models), and the standard bootstrap fails to capture this slower rate~\citep{kosorok2008,sen2015}.

The $m$-out-of-$n$ bootstrap draws $m < n$ samples without replacement, with $m = n^{2/3}$.
This restores consistency for shape-constrained estimators~\citep{leger2006}.
In our implementation, we use $m$-out-of-$n$ bootstrap for isotonic regression and standard bootstrap for logistic regression.


% ══════════════════════════════════════════════════════════════
\section{Results}
\label{sec:results}

\begin{table}[H]
\centering
\caption{Comparison of the three curve-fitting methods across all 18 models.
CV-Brier and CV-LogLoss are 5-fold cross-validated averages.
Doubling times and $R^2$ values come from exponential trend regression on per-model summary statistics vs.\ release date.}
\label{tab:results}
\begin{tabular}{@{}lcccccc@{}}
\toprule
\textbf{Method} & \textbf{CV-Brier} & \textbf{CV-LogLoss} & \textbf{$p_{50}$ Dbl} & \textbf{$p_{50}$ $R^2$} & \textbf{$G[T]$ Dbl} & \textbf{$G[T]$ $R^2$} \\
\midrule
Logistic (METR)     & 0.1212 & 0.376 & 121\,d & 0.916 & 128\,d & 0.916 \\
Isotonic            & 0.1115 & 0.368 & 118\,d & 0.811 & 130\,d & 0.921 \\
\bottomrule
\end{tabular}
\end{table}

\paragraph{Key takeaways.}
\begin{enumerate}
  \item \textbf{Fit quality:} Isotonic wins on both CV-Brier ($0.1115$ vs.\ $0.1212$) and CV-LogLoss ($0.368$ vs.\ $0.376$), outperforming logistic out of sample.
  \item \textbf{Metric stability:} $G[T]$ produces nearly identical doubling times ($128$--$130$\,d) and $R^2$ values ($0.916$--$0.921$) across both methods.
        Under $p_{50}$, the $R^2$ drops to $0.811$ for isotonic due to the step-function artifact, and individual-model $p_{50}$ values change dramatically between methods even though the trend slope is preserved.
  \item \textbf{Headline robustness:} The doubling time is approximately four months ($\sim$\!$128$--$130$ days) regardless of methodology.
  \item \textbf{Recommendation:} Use isotonic regression with $G[T]$ for the most defensible estimate. $G[T]$ captures the full fitted curve in a single number, making it robust to the choice of curve-fitting method.
\end{enumerate}


% ══════════════════════════════════════════════════════════════
\section{Interactive Companion}
\label{sec:interactive}

To enable detailed inspection of the results, we provide an interactive web-based viewer at:
\begin{center}
  \url{https://guyko81.github.io/metr-ablation/}
\end{center}
The viewer allows readers to:
\begin{itemize}
  \item Switch between both fitting methods and compare aggregate statistics side by side.
  \item Inspect per-model curve fits with bootstrap confidence bands for all 18 models.
  \item Compare $p_{50}$ and $G[T]$ trend charts with regression lines and confidence intervals.
  \item View detailed comparison tables of per-model metrics across methods.
\end{itemize}
All visualizations are interactive (Plotly-based) with hover tooltips for exact values.
The entire analysis is reproducible:
\begin{verbatim}
  pip install -r requirements.txt
  python run.py
\end{verbatim}


% ══════════════════════════════════════════════════════════════
\section{Discussion}
\label{sec:discussion}

\paragraph{The trend is real and robust.}
Our most important finding is negative in a positive sense: the methodology improvements do not overturn METR's headline result.
The doubling time of AI agent capabilities is approximately four months, and this holds across both logistic and isotonic curve fitting, and across both $p_{50}$ and $G[T]$ summary metrics.
What changes is the \emph{defensibility} of the estimate---a $G[T]$-based trend with $R^2 = 0.921$ is harder to dismiss than a $p_{50}$-based trend with $R^2 = 0.811$.

\paragraph{Policy implications.}
For decision-makers who rely on the METR doubling time to calibrate AI safety policy, our results are reassuring: the number is not an artifact of a particular curve-fitting choice.
The recommended methodology (isotonic + $G[T]$) provides a principled, well-calibrated, and stable estimate that should withstand methodological scrutiny.

\paragraph{Limitations and future work.}
\begin{itemize}
  \item \textbf{Trend regression:} We use ordinary least-squares on $\log(\text{metric})$ vs.\ release date, which assumes homoscedastic errors and perfect exponential growth.
        With only 18 data points, more sophisticated methods (Bayesian regression, distributional regression) are appealing but difficult to validate. We note this as future work.
  \item \textbf{Task weighting:} We use METR's default task weighting scheme. Sensitivity of the doubling time to alternative weightings deserves investigation.
  \item \textbf{Leave-one-model-out:} Testing stability of the trend when individual models are removed could further strengthen robustness claims.
\end{itemize}


% ══════════════════════════════════════════════════════════════
\section*{Acknowledgments}

We thank METR for making their benchmark data and methodology publicly available.
This work would not be possible without their systematic data collection effort.

\bibliographystyle{plainnat}
\begin{thebibliography}{9}

\bibitem[Barlow et~al.(1972)]{barlow1972}
R.~E. Barlow, D.~J. Bartholomew, J.~M. Bremner, and H.~D. Brunk.
\newblock \emph{Statistical Inference Under Order Restrictions}.
\newblock Wiley, 1972.

\bibitem[Kosorok(2008)]{kosorok2008}
M.~R. Kosorok.
\newblock \emph{Introduction to Empirical Processes and Semiparametric Inference}.
\newblock Springer, 2008.

\bibitem[Leger and MacGibbon(2006)]{leger2006}
C.~Leger and B.~MacGibbon.
\newblock On the bootstrap in cube root asymptotics.
\newblock \emph{Canadian Journal of Statistics}, 34(1):29--44, 2006.

\bibitem[METR(2025)]{metr2025}
METR.
\newblock Measuring AI agent time horizons, v1.1.
\newblock \url{https://metr.org/blog/2025-03-19-measuring-ai-agent-time-horizons/}, 2025.

\bibitem[Robertson et~al.(1988)]{robertson1988}
T.~Robertson, F.~T. Wright, and R.~L. Dykstra.
\newblock \emph{Order Restricted Statistical Inference}.
\newblock Wiley, 1988.

\bibitem[Sen and Xu(2015)]{sen2015}
B.~Sen and G.~Xu.
\newblock Model based bootstrap methods for interval censored data.
\newblock \emph{Computational Statistics \& Data Analysis}, 81:121--129, 2015.

\end{thebibliography}

\end{document}
